%%%%%%%%%%%%%%%%%%%%%%%%%%%%%%%%%%%%%%%%%
% Lab Report 5
% CS545_MachineLearning
% (2016SoE009)
%
%%%%%%%%%%%%%%%%%%%%%%%%%%%%%%%%%%%%%%%%%
%----------------------------------------------------------------------------------------
%	PACKAGES AND OTHER DOCUMENT CONFIGURATIONS
%----------------------------------------------------------------------------------------

\documentclass[12pt]{article}

\usepackage{fancyhdr} % Required for custom headers
\usepackage{lastpage} % Required to determine the last page for the footer
\usepackage{extramarks} % Required for headers and footers
\usepackage{graphicx} % Required to insert images
\usepackage{hyperref}
\usepackage{amsmath}
\usepackage{amsthm}
\usepackage{amssymb}
\usepackage{apacite}
\usepackage[english]{babel}
\usepackage{lipsum}
\usepackage{multirow}
\usepackage{booktabs}	%table rules
\usepackage[all]{nowidow}
\usepackage{threeparttable}
\usepackage{verbatim}
\usepackage{comment}
\usepackage{longtable}
\usepackage{etoolbox}
\setlength{\LTcapwidth}{=6.95in} %longtable caption width goes the full textwidth 



\usepackage{tikz}
\usetikzlibrary{shapes, arrows, positioning}
\tikzset{
    events/.style={ellipse, draw, align=right},
}

% paragraph indent
\newenvironment{myindentpar}[1]%
   {\begin{list}{}%
       {\setlength{\leftmargin}{#1}}%
           \item[]%
   }
     {\end{list}}

% Margins
\topmargin= -0.25in
\evensidemargin=0in
\oddsidemargin=0in
\textwidth=6.95in
\textheight=9.0in
\headsep=0.15in %distance between header and text

\linespread{1.2} % Line spacing


% Set up the header and footer
\pagestyle{fancy}
\lhead{} % Top left header
\chead{} % Top center header
\rhead{bmarron} % Top right header
\lfoot{\lastxmark} % Bottom left footer
\cfoot{} % Bottom center footer
\rfoot{Page\ \thepage\ of\ \pageref{LastPage}} % Bottom right footer
%\renewcommand\headrulewidth{0.4pt} % Size of the header rule
\renewcommand\footrulewidth{0.4pt} % Size of the footer rule

\setlength\parindent{0pt} % Removes all indentation from paragraphs
\setcounter{secnumdepth}{0} % Removes default section numbers

   
%----------------------------------------------------------------------------------------
%	NAME AND CLASS SECTION
%----------------------------------------------------------------------------------------

\newcommand{\hmwkClass}{Machine Learning (CS 545), Portland State University, Winter 2016} % Course/class
\newcommand{\hmwkTitle}{k-Means Clustering} % Assignment title
\newcommand{\hmwkAuthorName}{Bruce Marron} % Your name
\newcommand{\hmwkClassInstructor}{Mitchell} % Teacher/lecturer

%----------------------------------------------------------------------------------------
%	TITLE PAGE
%----------------------------------------------------------------------------------------

\title{
\vspace{2in}
\textmd{\textbf{\hmwkTitle}}\\
\textit{\normalsize \hmwkClass}
\vspace{3in}
}
\author{\textbf{\hmwkAuthorName}}
\date{08 March 2016} % Insert date here if you want it to appear below your name


%----------------------------------------
%	BEGIN DOC
%----------------------------------------
\begin{document}
\maketitle
\thispagestyle{empty}
\clearpage\maketitle

\subsection{Introduction}
The k-means algorithm is an algorithm for placing N data points $\mathbf{X} = \{ \mathbf{x}^{(n)}\} $ from an I-dimensional space into k clusters where each cluster is parameterized by a vector called its mean. Once the number of cluster has been defined and an initial parameterization of the means accomplished, the k-means algorithm proceeds as an iterative, two-step process. In the \textit{assignment step} each data point is assigned to the nearest mean (centroid). In the \textit{update step} means are adjusted to match the sample means of the data points for which they are responsible. The algorithm is run until the change in centroids is below some user-determined threshold (convergence).\\ 

This report documents the use of k-means clustering in two experiments to classify handwritten digits that have been extracted to normalized bitmaps given a feature vector in $ \mathbb{R}^{64}$. The "OptDigits" dataset was the source of raw data and was provided by the instructor. All data processing and evaluation tasks were done in Python 2.7.11 (Anaconda 2.4.1 (32-bit)) using the integrated development environment (IDE), "Spyder". The open source, "scikit-learn" machine learning package was used as a reference for code construction.\\

\subsection{Data Processing}
The raw data in OptDigits contains a training set of 3823 cases and a test set of 1797 cases. Each case is a vector of 64 features plus one class attribute. Feature values range from 0 - 16 and class attributes correspond to the natural number set, N = {0,1, 2,...,9}.  The data were imported and class attribute data were separated from case data for both the training dataset and the test dataset. This resulted in four, datasets: \verb!tr_d_X, tr_d_y, te_d_X, te_d_y!. The procedures and methods used for all data processing tasks are detailed in the script \verb!DataProcessing01.py!.

\subsection{Implementation of the k-means Algorithm} 
The implementation and actualization of the k-means algorithm is detailed in the following scripts,
\begin{verbatim}
DataProcessing02.py, DataProcessing03.py, DataProcessing04.py, DataProcessing05.py
DataProcessing06.py, DataProcessing07.py, DataProcessing08.py, DataProcessing09.py
DataProcessing10.py, DataProcessing11.py, DataProcessing12.py
\end{verbatim}


\subsection{Experiment 1: Results}
The results of Experiment 1 (k=10 clusters) are presented in Table~\ref{tab:a1} through Table~\ref{tab:a12}. Accuracy was determined as 20.14\% (=(176+3+2+1+173+2+5)/1797). As shown in Table~\ref{tab:a1}, the fourth random center cluster ('rcc4') had the minimum SSE and maximum entropy and was selected to generate the 'best centers' for Experiment 1. A confusion matrix was generated for each class to better detail classifier performance.

\begin{table} \centering 
  \caption{The SSE, sum-squared separation, and entropy for k-means classification (k=10) on the OptDigits training dataset.} 
  \label{tab:a1} 
\begin{tabular}{lccc} \\
[-1.8ex] \hline \hline \\[-1.8ex] 
 Initialization &&& \\ 
 (random cluster centers) & SSE & Sum-squared separation & Entropy\\
\hline \\[-1.8ex]
rcc1 & 2869491 & 166318 & 2.51909\\
rcc2 & 2716914 & 148726 & 2.94148\\
rcc3 & 2953240 & 161250 & 2.62772\\
rcc4 & 2596054 & 128680 & 3.09108\\
rcc5 & 2874411 & 159626 & 2.70705\\
\hline \\[-1.8ex]
\end{tabular} 
\end{table}


\begin{table} \centering 
  \caption{The assignment of cluster center labels from the best run of Experiment 1 (using rcc4) to the most frequent class they contain.} 
  \label{tab:a2} 
\begin{tabular}{lc} \\
[-1.8ex] \hline \hline \\[-1.8ex] 
 Cluster Center & Most Freq. Class \\ 
 \hline \\[-1.8ex]
0 & 0\\
7 & 1\\
7 & 2\\
8 & 3\\
5 & 4\\
9 & 5\\
6 & 6 \\
4 & 7\\
2 & 8\\
8 & 9\\
\hline \\[-1.8ex]
\end{tabular} 
\end{table}


\begin{table} \centering 
  \caption{The confusion matrix for Class = 0 using the 'best centers' k-means classifier from Experiment 1 as applied to the OptDigits test dataset.} 
  \label{tab:a3} 
\begin{tabular}{lccc} 
\\[-1.8ex]\hline 
\hline \\[-1.8ex] 
\multicolumn{2}{c}{} & \multicolumn{2}{c}{Predicted} \\
%\cline{3-4}
\multicolumn{2}{c}{} & 0 & !0\\
%\hline \\[-1.8ex]
\multirow{2}{*}{Actual} & 0 & 176 & 10\\
& !0 & 2 & 1609\\
\hline \\[-1.8ex]
\end{tabular} 
\end{table}

\begin{table} \centering 
  \caption{The confusion matrix for Class = 1 using the 'best centers' k-means classifier from Experiment 1 as applied to the OptDigits test dataset.} 
  \label{tab:a4} 
\begin{tabular}{lccc} 
\\[-1.8ex]\hline 
\hline \\[-1.8ex] 
\multicolumn{2}{c}{} & \multicolumn{2}{c}{Predicted} \\
%\cline{3-4}
\multicolumn{2}{c}{} & 1 & !1\\
%\hline \\[-1.8ex]
\multirow{2}{*}{Actual} & 1 & 0 & 0\\
& !1 & 182 & 1615\\
\hline \\[-1.8ex]
\end{tabular} 
\end{table}

\begin{table} \centering 
  \caption{The confusion matrix for Class = 2 using the 'best centers' k-means classifier from Experiment 1 as applied to the OptDigits test dataset.} 
  \label{tab:a5} 
\begin{tabular}{lccc} 
\\[-1.8ex]\hline 
\hline \\[-1.8ex] 
\multicolumn{2}{c}{} & \multicolumn{2}{c}{Predicted} \\
%\cline{3-4}
\multicolumn{2}{c}{} & 2 & !2\\
%\hline \\[-1.8ex]
\multirow{2}{*}{Actual} & 2 & 3 & 128\\
& !2 & 174 & 1492\\
\hline \\[-1.8ex]
\end{tabular} 
\end{table}

\begin{table} \centering 
  \caption{The confusion matrix for Class = 3 using the 'best centers' k-means classifier from Experiment 1 as applied to the OptDigits test dataset.} 
  \label{tab:a6} 
\begin{tabular}{lccc} 
\\[-1.8ex]\hline 
\hline \\[-1.8ex] 
\multicolumn{2}{c}{} & \multicolumn{2}{c}{Predicted} \\
%\cline{3-4}
\multicolumn{2}{c}{} & 3 & !3\\
%\hline \\[-1.8ex]
\multirow{2}{*}{Actual} & 3 & 2 & 131\\
& !3 & 181 & 1483\\
\hline \\[-1.8ex]
\end{tabular} 
\end{table}


\begin{table} \centering 
  \caption{The confusion matrix for Class = 4 using the 'best centers' k-means classifier from Experiment 1 as applied to the OptDigits test dataset.} 
  \label{tab:a7} 
\begin{tabular}{lccc} 
\\[-1.8ex]\hline 
\hline \\[-1.8ex] 
\multicolumn{2}{c}{} & \multicolumn{2}{c}{Predicted} \\
%\cline{3-4}
\multicolumn{2}{c}{} & 4 & !4\\
%\hline \\[-1.8ex]
\multirow{2}{*}{Actual} & 4 & 1 & 152\\
& !4 & 180 & 1464\\
\hline \\[-1.8ex]
\end{tabular} 
\end{table}

\begin{table} \centering 
  \caption{The confusion matrix for Class = 5 using the 'best centers' k-means classifier from Experiment 1 as applied to the OptDigits test dataset.} 
  \label{tab:a8} 
\begin{tabular}{lccc} 
\\[-1.8ex]\hline 
\hline \\[-1.8ex] 
\multicolumn{2}{c}{} & \multicolumn{2}{c}{Predicted} \\
%\cline{3-4}
\multicolumn{2}{c}{} & 5 & !5\\
%\hline \\[-1.8ex]
\multirow{2}{*}{Actual} & 5 & 0 & 187\\
& !5 & 182 & 1428\\
\hline \\[-1.8ex]
\end{tabular} 
\end{table}

\begin{table} \centering 
  \caption{The confusion matrix for Class = 6 using the 'best centers' k-means classifier from Experiment 1 as applied to the OptDigits test dataset.} 
  \label{tab:a9} 
\begin{tabular}{lccc} 
\\[-1.8ex]\hline 
\hline \\[-1.8ex] 
\multicolumn{2}{c}{} & \multicolumn{2}{c}{Predicted} \\
%\cline{3-4}
\multicolumn{2}{c}{} & 6 & !6\\
%\hline \\[-1.8ex]
\multirow{2}{*}{Actual} & 6 & 173 & 8\\
& !6 & 8 & 1608\\
\hline \\[-1.8ex]
\end{tabular} 
\end{table}

\begin{table} \centering 
  \caption{The confusion matrix for Class = 7 using the 'best centers' k-means classifier from Experiment 1 as applied to the OptDigits test dataset.} 
  \label{tab:a10} 
\begin{tabular}{lccc} 
\\[-1.8ex]\hline 
\hline \\[-1.8ex] 
\multicolumn{2}{c}{} & \multicolumn{2}{c}{Predicted} \\
%\cline{3-4}
\multicolumn{2}{c}{} & 7 & !7\\
%\hline \\[-1.8ex]
\multirow{2}{*}{Actual} & 7 & 2 & 336\\
& !7 & 177 & 1282\\
\hline \\[-1.8ex]
\end{tabular} 
\end{table}


\begin{table} \centering 
  \caption{The confusion matrix for Class = 8 using the 'best centers' k-means classifier from Experiment 1 as applied to the OptDigits test dataset.} 
  \label{tab:a11} 
\begin{tabular}{lccc} 
\\[-1.8ex]\hline 
\hline \\[-1.8ex] 
\multicolumn{2}{c}{} & \multicolumn{2}{c}{Predicted} \\
%\cline{3-4}
\multicolumn{2}{c}{} & 8 & !8\\
%\hline \\[-1.8ex]
\multirow{2}{*}{Actual} & 8 & 5 & 382\\
& !8 & 169 & 1295\\
\hline \\[-1.8ex]
\end{tabular} 
\end{table}


\begin{table} \centering 
  \caption{The confusion matrix for Class = 9 using the 'best centers' k-means classifier from Experiment 1 as applied to the OptDigits test dataset.} 
  \label{tab:a12} 
\begin{tabular}{lccc} 
\\[-1.8ex]\hline 
\hline \\[-1.8ex] 
\multicolumn{2}{c}{} & \multicolumn{2}{c}{Predicted} \\
%\cline{3-4}
\multicolumn{2}{c}{} & 9 & !9\\
%\hline \\[-1.8ex]
\multirow{2}{*}{Actual} & 9 & 0 & 155\\
& !9 & 180 & 1462\\
\hline \\[-1.8ex]
\end{tabular} 
\end{table}

\subsection{Experiment 2: Results}
The results of Experiment 2 (k=30 clusters) are presented in Table~\ref{tab:b1} through Table~\ref{tab:b12}. The accuracy was determined as 51.53\% (= 84+100+59+112+76+100+92+108+116+79)/1797). As shown in Table~\ref{tab:b1}, the fourth random center cluster ('rcc4') again had the minimum SSE and maximum entropy and was selected to generate the 'best centers' for Experiment 2. Again, a confusion matrix was generated for each class to better detail classifier performance.

\begin{table} \centering 
  \caption{The SSE, sum-squared separation, and entropy for k-means classification (k=30) on the OptDigits training dataset.} 
  \label{tab:b1} 
\begin{tabular}{lccc} \\
[-1.8ex] \hline \hline \\[-1.8ex] 
 Initialization &&& \\ 
 (random cluster centers) & SSE & Sum-squared separation & Entropy\\
\hline \\[-1.8ex]
rcc1 & 2209306 & 1698804 & 4.04929\\
rcc2 & 2438089 & 1272426 & 3.21328\\
rcc3 & 2121108 & 1713960 & 3.97819\\
rcc4 & 2012952 & 1692684 & 4.42242\\
rcc5 & 2585148 & 1267474 & 3.12321\\
\hline \\[-1.8ex]
\end{tabular} 
\end{table}


\begin{table} \centering 
  \caption{The assignment of cluster center labels from the best run of Experiment 2 (using rcc4) to the most frequent class they contain.} 
  \label{tab:b2} 
\begin{tabular}{lc} \\
[-1.8ex] \hline \hline \\[-1.8ex] 
 Cluster Center & Most Freq. Class \\ 
 \hline \\[-1.8ex]
13 & 0\\
9 & 1\\
7 & 2\\
27 & 3\\
19 & 4\\
0 & 5\\
26 & 6 \\
23 & 7\\
10 & 8\\
20 & 9\\
\hline \\[-1.8ex]
\end{tabular} 
\end{table}

\begin{table} \centering 
  \caption{The confusion matrix for Class = 0 using the 'best centers' k-means classifier from Experiment 2 as applied to the OptDigits test dataset.} 
  \label{tab:b3} 
\begin{tabular}{lccc} 
\\[-1.8ex]\hline 
\hline \\[-1.8ex] 
\multicolumn{2}{c}{} & \multicolumn{2}{c}{Predicted} \\
%\cline{3-4}
\multicolumn{2}{c}{} & 0 & !0\\
%\hline \\[-1.8ex]
\multirow{2}{*}{Actual} & 0 & 84 & 0\\
& !0 & 94 & 1619\\
\hline \\[-1.8ex]
\end{tabular} 
\end{table}

\begin{table} \centering 
  \caption{The confusion matrix for Class = 1 using the 'best centers' k-means classifier from Experiment 2 as applied to the OptDigits test dataset.} 
  \label{tab:b4} 
\begin{tabular}{lccc} 
\\[-1.8ex]\hline 
\hline \\[-1.8ex] 
\multicolumn{2}{c}{} & \multicolumn{2}{c}{Predicted} \\
%\cline{3-4}
\multicolumn{2}{c}{} & 1 & !1\\
%\hline \\[-1.8ex]
\multirow{2}{*}{Actual} & 1 & 100 & 23\\
& !1 & 82 & 1592\\
\hline \\[-1.8ex]
\end{tabular} 
\end{table}

\begin{table} \centering 
  \caption{The confusion matrix for Class = 2 using the 'best centers' k-means classifier from Experiment 2 as applied to the OptDigits test dataset.} 
  \label{tab:b5} 
\begin{tabular}{lccc} 
\\[-1.8ex]\hline 
\hline \\[-1.8ex] 
\multicolumn{2}{c}{} & \multicolumn{2}{c}{Predicted} \\
%\cline{3-4}
\multicolumn{2}{c}{} & 2 & !2\\
%\hline \\[-1.8ex]
\multirow{2}{*}{Actual} & 2 & 59 & 20\\
& !2 & 118 & 1600\\
\hline \\[-1.8ex]
\end{tabular} 
\end{table}

\begin{table} \centering 
  \caption{The confusion matrix for Class = 3 using the 'best centers' k-means classifier from Experiment 2 as applied to the OptDigits test dataset.} 
  \label{tab:b6} 
\begin{tabular}{lccc} 
\\[-1.8ex]\hline 
\hline \\[-1.8ex] 
\multicolumn{2}{c}{} & \multicolumn{2}{c}{Predicted} \\
%\cline{3-4}
\multicolumn{2}{c}{} & 3 & !3\\
%\hline \\[-1.8ex]
\multirow{2}{*}{Actual} & 3 & 112 & 3\\
& !3 & 71 & 1611\\
\hline \\[-1.8ex]
\end{tabular} 
\end{table}

\clearpage	%add if "too many unprocessed floats"error
\begin{table} \centering 
  \caption{The confusion matrix for Class = 4 using the 'best centers' k-means classifier from Experiment 2 as applied to the OptDigits test dataset.} 
  \label{tab:b7} 
\begin{tabular}{lccc} 
\\[-1.8ex]\hline 
\hline \\[-1.8ex] 
\multicolumn{2}{c}{} & \multicolumn{2}{c}{Predicted} \\
%\cline{3-4}
\multicolumn{2}{c}{} & 4 & !4\\
%\hline \\[-1.8ex]
\multirow{2}{*}{Actual} & 4 & 76 & 2\\
& !4 & 105 & 1614\\
\hline \\[-1.8ex]
\end{tabular} 
\end{table}

\begin{table}[!htbp] \centering 
  \caption{The confusion matrix for Class = 5 using the 'best centers' k-means classifier from Experiment 2 as applied to the OptDigits test dataset.} 
  \label{tab:b8} 
\begin{tabular}{lccc} 
\\[-1.8ex]\hline 
\hline \\[-1.8ex] 
\multicolumn{2}{c}{} & \multicolumn{2}{c}{Predicted} \\
%\cline{3-4}
\multicolumn{2}{c}{} & 5 & !5\\
%\hline \\[-1.8ex]
\multirow{2}{*}{Actual} & 5 & 100 & 3\\
& !5 & 82 & 1612\\
\hline \\[-1.8ex]
\end{tabular} 
\end{table}

\begin{table}[!htbp] \centering 
  \caption{The confusion matrix for Class = 6 using the 'best centers' k-means classifier from Experiment 2 as applied to the OptDigits test dataset.} 
  \label{tab:b9} 
\begin{tabular}{lccc} 
\\[-1.8ex]\hline 
\hline \\[-1.8ex] 
\multicolumn{2}{c}{} & \multicolumn{2}{c}{Predicted} \\
%\cline{3-4}
\multicolumn{2}{c}{} & 6 & !6\\
%\hline \\[-1.8ex]
\multirow{2}{*}{Actual} & 6 & 92 & 0\\
& !6 & 89 & 1616\\
\hline \\[-1.8ex]
\end{tabular} 
\end{table}

\begin{table}[!htbp] \centering 
  \caption{The confusion matrix for Class = 7 using the 'best centers' k-means classifier from Experiment 2 as applied to the OptDigits test dataset.} 
  \label{tab:b10} 
\begin{tabular}{lccc} 
\\[-1.8ex]\hline 
\hline \\[-1.8ex] 
\multicolumn{2}{c}{} & \multicolumn{2}{c}{Predicted} \\
%\cline{3-4}
\multicolumn{2}{c}{} & 7 & !7\\
%\hline \\[-1.8ex]
\multirow{2}{*}{Actual} & 7 & 108 & 12\\
& !7 & 71 & 1606\\
\hline \\[-1.8ex]
\end{tabular} 
\end{table}


\begin{table}[!htbp] \centering 
  \caption{The confusion matrix for Class = 8 using the 'best centers' k-means classifier from Experiment 2 as applied to the OptDigits test dataset.} 
  \label{tab:b11} 
\begin{tabular}{lccc} 
\\[-1.8ex]\hline 
\hline \\[-1.8ex] 
\multicolumn{2}{c}{} & \multicolumn{2}{c}{Predicted} \\
%\cline{3-4}
\multicolumn{2}{c}{} & 8 & !8\\
%\hline \\[-1.8ex]
\multirow{2}{*}{Actual} & 8 & 116 & 14\\
& !8 & 58 & 1609\\
\hline \\[-1.8ex]
\end{tabular} 
\end{table}


\begin{table}[!htbp] \centering 
  \caption{The confusion matrix for Class = 9 using the 'best centers' k-means classifier from Experiment 2 as applied to the OptDigits test dataset.} 
  \label{tab:b12} 
\begin{tabular}{lccc} 
\\[-1.8ex]\hline 
\hline \\[-1.8ex] 
\multicolumn{2}{c}{} & \multicolumn{2}{c}{Predicted} \\
%\cline{3-4}
\multicolumn{2}{c}{} & 9 & !9\\
%\hline \\[-1.8ex]
\multirow{2}{*}{Actual} & 9 & 79 & 21\\
& !9 & 101 & 1596\\
\hline \\[-1.8ex]
\end{tabular} 
\end{table}

\subsection{Discussion}
The accuracy certainly improved with k=30 clusters although it is still far short of the results reported by the authors of the OptDigit dataset (~90\%). This is most likely to the random selection for all starting clusters rather than a pre-selected set of initial clusters. Interestingly, the top two digits correctly classified differed between Experiment 1 and Experiment 2. In the first experiment, digits "0"and "6" were recognized with very high frequency (176/178 and 173/181, respectively) while in the second experiment, "8" and "3" were recognized with medium frequency (116/174 and 112/183, respectively).\\

All of the 'best center' clusters were translated to 8x8 bitmaps (.pmg files; Exp1Graphics and Exp2Graphics folders). Visualizing these images shows that many do resemble the digits being classified, as for example, 'best center' cluster 0 from Experiment 1 (Figure~\ref{fig:a1}) and "best center" cluster 8 from Experiment 2 (Figure~\ref{fig:a2} ). Others appear strange, as for example, "best center" cluster 5 from Experiment 1 (Figure~\ref{fig:a3} )and "best center" cluster 7 from Experiment 2 (Figure~\ref{fig:a4} ).

\begin{figure}
   	\begin{center}
   		\includegraphics[scale=0.5]{Screenshot_Exp1dig0}
   		\caption{'Best center' cluster 0 from Experiment 1. This centroid had good success in correctly identifying the digit, "0" from the OptDigits training dataset. }
   		\label{fig:a1}
   	\end{center}
   \end{figure}
   
\begin{figure}
   	\begin{center}
   		\includegraphics[scale=0.5]{Screenshot_Exp2dig10}
   		\caption{'Best center' cluster 10 from Experiment 2. This centroid had good success in correctly identifying the digit, "8" from the OptDigits training dataset.}
   		\label{fig:a2}
   	\end{center}
   \end{figure} 
   
 
\begin{figure}
   	\begin{center}
   		\includegraphics[scale=0.5]{Screenshot_Exp1dig5}
   		\caption{'Best center' cluster 5 from Experiment 1. This centroid had poor success in correctly identifying the digit, "4" from the OptDigits training dataset. }
   		\label{fig:a3}
   	\end{center}
   \end{figure}   

\begin{figure}  
	\begin{center}
  	 \includegraphics[scale=0.5]{Screenshot_Exp2dig7}
   		\caption{'Best center' cluster 7 from Experiment 2. This centroid had poor success in correctly identifying the digit, "2" from the OptDigits training dataset.}
   		\label{fig:a4}
   	\end{center}
   \end{figure}


      

\subsection{Conclusions}
Despite its limitations (sensitivity to initial cluster selection; sensitivity to outliers; lack of rigorous methods for selection of the number of clusters; local optimization) this algorithm appears to be very useful for general sorting purposes, especially on new datasets that lack well-defined priors.


%\renewcommand{\refname}{\normalfont\selectfont\small \textbf{References}} 
%\bibliographystyle{/usr/local/share/texmf/tex/latex/apacite/apacite}
%\bibliography{/home/bmarron//Desktop/BibTex/2016SoE009}

\end{document}
